\documentclass[11pt,letterpaper,landscape]{article}

\usepackage[
  top=0.55cm,
  bottom=0.55cm,
  left=0.65cm,
  right=0.65cm
]{geometry}
\usepackage{xcolor}
\usepackage{fontspec}
\usepackage{tcolorbox}
\tcbuselibrary{skins}
\usepackage{tabularx}
\usepackage{listings}
\usepackage{amsmath}

\setmainfont{Libertinus Sans}
\setmonofont{Libertinus Mono}
\pagestyle{empty}
\setlength{\parindent}{0pt}
\setlength{\tabcolsep}{4pt}
\renewcommand{\arraystretch}{1.12}
\newcolumntype{Y}{>{\raggedright\arraybackslash}X}

\definecolor{tealbox}{HTML}{079A9A}
\definecolor{orangeTitle}{HTML}{FF5A00}

\lstdefinestyle{pyinline}{
  basicstyle=\ttfamily\footnotesize,
  columns=fullflexible,
  keepspaces=true,
  breaklines=true,
  showstringspaces=false
}
\lstset{style=pyinline}

\newcommand{\code}[1]{\texttt{\detokenize{#1}}}
\newcommand{\entry}[2]{\textbf{\code{#1}} & #2\\}
\newcommand{\smallentry}[2]{\code{#1} & #2\\}
\newcommand{\colw}{0.318\textwidth}

\newtcolorbox{cheatbox}[1]{
  enhanced,
  sharp corners,
  boxrule=0pt,
  colback=white,
  left=2.5mm,
  right=2.5mm,
  top=1.2mm,
  bottom=1.2mm,
  before skip=2.5mm,
  after skip=2.5mm,
  title=#1,
  coltitle=white,
  fonttitle=\bfseries\large,
  colbacktitle=tealbox,
  attach boxed title to top center={yshift=-1mm},
  boxed title style={sharp corners, boxrule=0pt, colback=tealbox, left=3mm, right=3mm},
}

\newcommand{\cheattitle}{
  \begin{center}
    {\fontsize{23}{25}\selectfont\bfseries\color{orangeTitle} Python para Solemne 2 - Cheat Sheet}
  \end{center}
  \vspace{-3mm}
}

\begin{document}
\cheattitle

\noindent
\begin{minipage}[t]{\colw}
\vspace{0pt}
\begin{cheatbox}{Python: sintaxis básica}
\begin{tabularx}{\linewidth}{@{}lY@{}}
\textbf{\texttt{\# comentario}} & texto que Python ignora\\
\entry{x = 3.5}{asignar un valor}
\entry{a + b, a - b}{suma y resta}
\entry{a * b, a / b}{multiplicación y división}
\entry{a ** 2}{potencia}
\entry{(a + b) / 2}{ordenar cálculos}
\entry{print(x)}{mostrar resultado}
\entry{print(f"{x:.3f}")}{imprimir con 3 decimales}
\end{tabularx}
\end{cheatbox}

\begin{cheatbox}{Funciones}
\begin{lstlisting}
def nombre_funcion(t):
    resultado = ...
    return resultado

y = nombre_funcion(t)
\end{lstlisting}
Use funciones para separar el modelo físico del análisis numérico.
\end{cheatbox}

\begin{cheatbox}{Condicionales}
\begin{lstlisting}
if condicion:
    clasificacion = "alta"
else:
    clasificacion = "baja"
\end{lstlisting}
\begin{tabularx}{\linewidth}{@{}lY@{}}
\entry{>=, <=}{mayor/menor o igual}
\entry{>, <}{mayor o menor}
\entry{==}{igualdad}
\entry{and}{dos condiciones a la vez}
\end{tabularx}
\end{cheatbox}
\end{minipage}
\hfill
\begin{minipage}[t]{\colw}
\vspace{0pt}
\begin{cheatbox}{Listas e índices}
\begin{tabularx}{\linewidth}{@{}lY@{}}
\entry{lista = []}{crear lista vacía}
\entry{lista.append(x)}{agregar \code{x} al final}
\entry{L[i]}{elemento de \code{L} en índice \code{i}}
\entry{d[i]}{elemento de \code{d} en índice \code{i}}
\entry{len(d)}{cantidad de elementos}
\entry{range(n)}{índices desde \code{0} hasta \code{n-1}}
\end{tabularx}
\begin{lstlisting}
valores = []
for i in range(len(d)):
    calculo_F = L[i] / (4*np.pi*d[i]**2)
    valores.append(calculo_F)
\end{lstlisting}
\end{cheatbox}

\begin{cheatbox}{Importar módulos}
\begin{lstlisting}
import numpy as np
import matplotlib.pyplot as plt
\end{lstlisting}
\begin{tabularx}{\linewidth}{@{}lY@{}}
\entry{np}{alias usual de NumPy}
\entry{plt}{alias usual de Matplotlib}
\end{tabularx}
\end{cheatbox}

\begin{cheatbox}{NumPy: crear arreglos}
\begin{tabularx}{\linewidth}{@{}lY@{}}
\entry{np.array([1,2,3])}{crear arreglo desde una lista}
\entry{np.linspace(0,12,300)}{300 valores entre 0 y 12}
\entry{np.zeros_like(t)}{ceros con forma de \code{t}}
\entry{np.ones_like(t)}{unos con forma de \code{t}}
\entry{c*np.ones_like(t)}{arreglo constante igual a \code{c}}
\entry{np.ones(n)}{arreglo de \code{n} unos}
\end{tabularx}
\end{cheatbox}
\end{minipage}
\hfill
\begin{minipage}[t]{\colw}
\vspace{0pt}
\begin{cheatbox}{NumPy: funciones matemáticas}
\begin{tabularx}{\linewidth}{@{}lY@{}}
\entry{np.pi}{valor de $\pi$}
\entry{np.sin(x)}{seno de \code{x} en radianes}
\entry{np.cos(x)}{coseno de \code{x} en radianes}
\entry{np.deg2rad(theta)}{convertir grados a radianes}
\entry{t**2}{cuadrado elemento a elemento}
\entry{a*t + b}{operación vectorizada}
\end{tabularx}
\end{cheatbox}

\begin{cheatbox}{NumPy: análisis de datos}
\begin{tabularx}{\linewidth}{@{}lY@{}}
\entry{np.max(B)}{máximo del arreglo}
\entry{np.min(B)}{mínimo del arreglo}
\entry{np.mean(B)}{promedio}
\entry{np.sum(B > prom)}{cuenta valores que cumplen condición}
\entry{np.argmax(v)}{índice del máximo}
\entry{B > prom}{comparación elemento a elemento}
\end{tabularx}
\end{cheatbox}

\begin{cheatbox}{Modelos vectorizados}
\begin{lstlisting}
t = np.linspace(0, 12, 300)
B = 1 + a1*np.sin(2*np.pi*t/p)

theta = np.deg2rad(theta_deg)
x = v0*np.cos(theta)*t
y = y0 + v0*np.sin(theta)*t - 0.5*g*t**2
\end{lstlisting}
Si \code{t} es un arreglo, las operaciones producen arreglos.
\end{cheatbox}
\end{minipage}

\newpage
\cheattitle

\noindent
\begin{minipage}[t]{\colw}
\vspace{0pt}
\begin{cheatbox}{Matplotlib: crear figura}
\begin{tabularx}{\linewidth}{@{}lY@{}}
\entry{plt.figure(figsize=(9,5))}{crear figura}
\entry{plt.show()}{mostrar figura}
\entry{plt.tight_layout()}{ajustar espacios}
\end{tabularx}
\end{cheatbox}

\begin{cheatbox}{Matplotlib: graficar}
\begin{tabularx}{\linewidth}{@{}lY@{}}
\entry{plt.plot(t,B)}{gráfico de línea}
\entry{plt.scatter(d,F)}{gráfico de dispersión}
\entry{plt.axhline(prom)}{línea horizontal}
\end{tabularx}
\end{cheatbox}

\begin{cheatbox}{Matplotlib: presentación}
\begin{tabularx}{\linewidth}{@{}lY@{}}
\entry{plt.title("...")}{título}
\entry{plt.xlabel("...")}{nombre eje x}
\entry{plt.ylabel("...")}{nombre eje y}
\entry{plt.grid(True)}{activar grilla}
\entry{plt.legend()}{mostrar leyenda}
\end{tabularx}
\end{cheatbox}
\end{minipage}
\hfill
\begin{minipage}[t]{\colw}
\vspace{0pt}
\begin{cheatbox}{Subgráficos: crear grilla}
\begin{lstlisting}
fig, axs = plt.subplots(3, 2, figsize=(10, 9))
\end{lstlisting}
\begin{tabularx}{\linewidth}{@{}lY@{}}
\entry{3, 2}{3 filas y 2 columnas}
\entry{axs[fila,col]}{seleccionar un panel}
\end{tabularx}
\end{cheatbox}

\begin{cheatbox}{Subgráficos: usar paneles}
\begin{lstlisting}
axs[0, 0].plot(t, y)
axs[0, 0].set_title("Posición vertical")
axs[0, 0].set_xlabel("t [s]")
axs[0, 0].set_ylabel("y(t) [m]")
axs[0, 0].grid(True)
\end{lstlisting}
\end{cheatbox}

\begin{cheatbox}{Subgráficos: posiciones}
\begin{tabularx}{\linewidth}{@{}lY@{}}
\entry{axs[0,0]}{fila 1, columna 1}
\entry{axs[1,0]}{fila 2, columna 1}
\entry{axs[2,0]}{fila 3, columna 1}
\entry{axs[0,1]}{fila 1, columna 2}
\entry{axs[1,1]}{fila 2, columna 2}
\entry{axs[2,1]}{fila 3, columna 2}
\end{tabularx}
\end{cheatbox}
\end{minipage}
\hfill
\begin{minipage}[t]{\colw}
\vspace{0pt}
\begin{cheatbox}{Errores típicos: ángulos}
\begin{tabularx}{\linewidth}{@{}lY@{}}
\smallentry{sin(55)}{mal si 55 está en grados; use \code{np.deg2rad}}
\smallentry{np.sin(theta)}{\code{theta} debe estar en radianes}
\end{tabularx}
\end{cheatbox}

\begin{cheatbox}{Errores típicos: listas}
\begin{tabularx}{\linewidth}{@{}lY@{}}
\smallentry{lista.append}{requiere paréntesis: \code{lista.append(x)}}
\smallentry{L[i], d[i]}{use el mismo índice para datos relacionados}
\smallentry{range(len(d))}{recorre todos los índices de \code{d}}
\end{tabularx}
\end{cheatbox}

\begin{cheatbox}{Errores típicos: arreglos}
\begin{tabularx}{\linewidth}{@{}lY@{}}
\smallentry{vx = constante}{para graficar, conviene usar \code{ones_like}}
\smallentry{axs[0].plot}{en grilla 2D use \code{axs[0,0].plot}}
\smallentry{if B > prom}{\code{B > prom} es un arreglo; use \code{np.sum}}
\end{tabularx}
\end{cheatbox}
\end{minipage}

\end{document}
